\chapter{TRÌNH BÀY GIẢI PHÁP}\label{chap:giaiphap}

Phần trình bày dùng bộ slide \code{ScholarLens\_edited.pptx} (15 trang, có ghi chú người nói), sắp xếp theo đúng
các đề mục của báo cáo. Phần demo chạy trực tiếp trên web app với 4 PDF trong \code{demo\_pdfs/}; cả bốn bài chưa
có trong kho trước buổi demo.

\section{Kịch bản demo}

\begin{enumerate}[label=(\arabic*)]
\item \textbf{Chuẩn bị:} trên Kaggle chạy notebook \code{kaggle\_llm\_server.ipynb}; trên máy demo chạy
  \code{python -m scholarlens.app}, mở \code{http://127.0.0.1:7860}, dán URL trycloudflare và API key vào dòng
  ``LLM'' ở cuối thanh bên. Chấm xanh là sẵn sàng.
\item \textbf{Thêm một bài:} kéo file GraphRAG (\code{2404.16130.pdf}) vào cửa sổ. Tiến trình hiện ngay trong khung
  chat: SHA-256 $\to$ đọc PDF $\to$ tra arXiv API $\to$ chia đoạn $\to$ embedding $\to$ đồ thị. Kết thúc, dòng
  ``Đang hỏi về: \ldots'' xuất hiện phía trên ô nhập.
\item \textbf{Hỏi về bài:} ``Tóm tắt bài này'', ``Phương pháp chính và đóng góp là gì?''. Bấm một trích dẫn [n] để
  mở danh sách nguồn và tô sáng đúng bài đó.
\item \textbf{Hỏi bài liên quan:} ``Các bài báo liên quan đến bài này?'' cho danh sách kèm lý do (trích dẫn, cùng
  khái niệm, tương đồng nội dung). Bấm nút Đồ thị dưới câu trả lời, rồi Mở rộng để xem toàn màn hình và bấm vào một
  bài để xem chi tiết.
\item \textbf{Thêm trùng:} kéo lại cùng file; hệ thống báo ``đã có --- trùng nội dung file (SHA-256)'' và không
  embedding lại.
\item \textbf{Thêm hàng loạt:} chọn cùng lúc 3 PDF còn lại (RAG Survey, LightRAG, Agentic RAG); mỗi file hiện trạng
  thái riêng, rồi hỏi ``So sánh các bài vừa thêm''.
\item \textbf{Lịch sử và trình chiếu:} mở lại một hội thoại cũ trong thanh bên để cho thấy câu trả lời, nguồn và đồ
  thị được lưu lại; thu gọn thanh bên để khung chat chiếm toàn màn hình khi trình chiếu.
\end{enumerate}

\section{Kết luận}

ScholarLens đáp ứng đủ các yêu cầu của đề bài:
\begin{itemize}
\item Thu thập bài báo từ arXiv theo keyword, theo chuyên mục và theo mã bài, cùng nguồn snapshot Kaggle.
\item Tổ chức bốn kho dữ liệu (SQLite, ChromaDB, BM25, đồ thị) dựng lại được từ SQLite, với kiểm tra trùng ba lớp.
\item Tìm kiếm lai BM25 + vector đa ngôn ngữ gộp bằng RRF, hỏi được bằng tiếng Việt trên kho tiếng Anh.
\item Chatbot Graph RAG hai chế độ, trả lời có trích dẫn và gợi ý bài liên quan có giải thích.
\item LLM mã nguồn mở tự host bằng vLLM trên GPU miễn phí, chịu được việc mất kết nối.
\end{itemize}
Thao tác thêm PDF được tự động hóa trọn vẹn, kể cả kiểm tra trùng, và hệ thống vẫn chạy được khi không có LLM.

\section{Hướng phát triển}

\begin{table}[H]
\centering\small\setstretch{1.15}
\begin{tabularx}{\linewidth}{@{}P{0.36\linewidth}L@{}}
\toprule
\textbf{Hạn chế hiện tại} & \textbf{Hướng cải tiến}\\
\midrule
Đọc PDF bằng quy tắc & Dùng GROBID hoặc Nougat để lấy cấu trúc, bảng, công thức\\
Cạnh trích dẫn chỉ qua mã arXiv & Khớp tiêu đề tài liệu tham khảo, dùng Semantic Scholar API\\
Khái niệm dạng n-gram & Trích thực thể và quan hệ bằng LLM, tóm tắt cộng đồng như GraphRAG~\cite{graphrag} đầy đủ\\
Chưa có bước xếp hạng lại & Thêm cross-encoder (ví dụ bge-reranker) cho top-20 kết quả\\
Thu thập chạy thủ công & Theo dõi chủ đề định kỳ, chỉ lấy bài mới kể từ lần trước rồi tự nạp\\
Đánh giá known-item còn dễ & Xây bộ câu hỏi tiếng Việt có gán nhãn, đánh giá câu trả lời bằng LLM\\
\bottomrule
\end{tabularx}
\caption{Hạn chế và hướng phát triển}\label{tab:huongphattrien}
\end{table}
